\chapter{Introduction}\label{intro}
Money laundering is the criminal process of disguising or concealing the origin, identity, and ownership of illegally acquired funds.
Its detection has become a foundational pillar of modern economic security.
The rapid digitization of financial services has opened new possibilities for fraudulent activities.
Undetected ML leads to massive economic losses for individuals and institutions, undermines the public trust in the financial system and frequently serves to finance criminal organizations.
The scale of this phenomenon is staggering: the United Nations Office on Drugs and Crime estimates that between 2\% and 5\% of global GDP is laundered annually, with less than 1\% of these illicit financial flows being successfully detected~\cite{unodc2011estimating,europol2016does}.
According to the Canadian Anti Fraud Centre, 704 million dollars were lost because of frauds in 2025 only, confirming the raising trend of the previous years \cite{cafc_home}. 
Consequently, developing robust mechanisms to identify and stop it is critically important to preserving both market integrity and social stability.

The goal of this dissertation is to evaluate the effectiveness of Graph Neural Network (GNN) based models on detecting money laundering patterns. Both precision and explainability will be taken into account and we will propose two different tasks; a binary classification to label transactions as licit or illicit, and a multiclass classification to detect the specific type of laundering strategy employed which, in case of succes, will enhance the explainability of the model. 

In particular we will focus on the Group Aware Graph Neural Network (GAGNN) architecture proposed in \cite{mainpaper}. GAGNN is a GNN-based model that leverages a multi-component loss function to predict if a transaction is licit or illicit. There are a total of 3 losses, a node loss, an edge loss and a group loss. The group one is the novelty of the proposal since it introduces the concept of community by grouping nodes that are suspected to be doing money laundering toghether. More details on the architecture will be provided in Section \ref{sec:gagnn}. 
We will investigate the following inquiries;
\begin{itemize}
    \item whether the GAGNN architecture outperforms other baseline models on the binary task;
    \item whether the GAGNN architecture could be improved by enhancing Graph Isomorphism Network (GIN) instead of Graph Attention Network (GAT);
    \item whether the GAGNN architecture could perform well also on the multiclass task.
\end{itemize}

\section{Structure of the Dissertation}
The dissertation is organized as follows.

\textbf{Chapter~\ref{background} --- Background.}
This chapter lays the theoretical foundations of the work. It begins by introducing the mechanics of money laundering, covering smurfing, fraud rings, and the eight structural laundering patterns, together with an overview of the applicable regulatory framework and the ethical issues that arise when deploying machine learning models in the financial domain. The chapter then motivates the graph-based formulation of the problem, introduces the field of Graph Machine Learning, and formalizes the message-passing framework underpinning GNNs. Two specific architectures are described in detail: Graph Attention Networks (GAT, Section~\ref{sec:gat}) and Graph Isomorphism Networks (GIN, Section~\ref{sec:gin}), followed by a discussion of their theoretical limitations (Section \ref{sec:mp_limits}).

\textbf{Chapter~\ref{approaches} --- Methodology.}
This chapter presents the proposed approach. It describes the Group-Aware Graph Neural Network (GAGNN) architecture, detailing each component of its pipeline: the community-centric encoder, the eMRF layer, the node-level classifier, the group representation layer, and the multi-task loss.
Then it proceeds by providing the description of the \textit{IBM AMLSim Small} dataset used in this work, toghether with his features distributions.
Finally, it covers the experimental setup, including data preparation, the train/validation/test split strategy, subgraph sampling, and the two-phase model selection procedure used to tune both the loss function hyperparameters and the architectural ones.
\textbf{Chapter~\ref{results} --- Results.}
This chapter presents and discusses the experimental results. After reporting the outcomes of the model selection phase, the binary and multiclass classification results of the four models (GAGNN-GAT, GAGNN-GIN, baseline GAT, and baseline GIN) are compared using accuracy, precision, recall, F1-Score, and ROC-AUC as metrics. The chapter concludes with a discussion of the limitations encountered during the study.

\textbf{Chapter~\ref{sec:concl} --- Conclusions.}
The final chapter revisits the three research inquiries formulated above, interprets the results in light of the theoretical properties of the models, and outlines the most promising directions for future work.
